\section{General homogeneous covariance in dimension two}
\label{S:planar-holder}

\myCmts{Candidate argument, pending independent review. This section gives
complete endpoint constructions and a complete planar transfer. It does not
use the preceding candidate proposition as an accepted input. Earlier
accepted-scope statements and review confidence remain unchanged.}

Throughout this section $d=2$. Write $\Gamma$ for the spatial covariance
as a nonnegative, even, positive-definite tempered Radon measure. For a
locally integrable covariance function, this means
$\Gamma(\ud x)=\gamma(x)\ud x$. The measure interpretation uses the Gaussian
approximation framework of
Chen~\cite{chen:author2026:time-dependency}*{pp.~2--3, 13} and
Chen--Hu~\cite{chen.hu:25:hyperbolic}*{pp.~2145--2147, 2151--2153}.
All products of covariance factors below are taken after Gaussian smoothing;
no product or fractional power of distributions is used. Homogeneity means
\begin{equation}\label{E:planar-measure-scaling}
  \Gamma(rA)=r^{2-\alpha}\Gamma(A),\qquad r>0,
\end{equation}
for Borel sets $A$. This is equivalent to the covariance homogeneity in
Hypothesis~\ref{H:homog}. If $\Gamma=0$, the noise vanishes and $u=1$.
We first treat $\Gamma\ne0$ and $1<\alpha<2$.

\begin{lemma}[Planar endpoint covariances]\label{L:planar-endpoints}
  Suppose that $\Gamma$ satisfies the preceding assumptions, and let
  $\alpha<\beta<2$. Define the finite angular measure
  \[
    \sigma(E)=(2-\alpha)\Gamma
    \{x:0<|x|\le1,\ x/|x|\in E\},\qquad
    S=\sigma(\mathbb{S}^1).
  \]
  On $\R^2\setminus\{0\}$ define the weighted measures
  \begin{equation}\label{E:planar-endpoints}
    \Gamma_{\rm low}(\ud x)=|x|^{\alpha-1}\Gamma(\ud x),\qquad
    \Gamma_{\rm st}(\ud x)=|x|^{\alpha-\beta}\Gamma(\ud x),
  \end{equation}
  with zero mass at the origin. Both are nonnegative, even,
  positive-definite tempered Radon measures. Their measure scaling
  exponents are $1$ and $2-\beta$, respectively. For every $h>0$,
  $\epsilon>0$, $\eta\ge0$, and $z\in\R^2$,
  \begin{equation}\label{E:planar-line-convolution}
    \int G_{h,\eta}(\ud y)(p_\epsilon*\Gamma_{\rm low})(y+z)
    \le\frac S4.
  \end{equation}
  If $\vartheta=(\beta-\alpha)/(\beta-1)$, then $0<\vartheta<1$ and
  \begin{equation}\label{E:planar-gaussian-holder}
    (p_\epsilon*\Gamma)(z)
    \le(p_\epsilon*\Gamma_{\rm low})(z)^\vartheta
    (p_\epsilon*\Gamma_{\rm st})(z)^{1-\vartheta}.
  \end{equation}
  No angular density or isotropy is assumed.
\end{lemma}

\begin{proof}
  Since $2-\alpha>0$, scaling and local finiteness force
  $\Gamma(\{0\})=0$. The unit-ball mass is finite, so $S<\infty$.
  It is positive: otherwise scaling would make every ball have zero
  mass, contradicting $\Gamma\ne0$. Evenness of $\Gamma$ implies
  evenness of $\sigma$. For $0<a<b$ and Borel $E\subset\mathbb{S}^1$,
  \[
    \Gamma\{x:a<|x|\le b,\ x/|x|\in E\}
    =\frac{b^{2-\alpha}-a^{2-\alpha}}{2-\alpha}\sigma(E)
    =\int_a^b r^{1-\alpha}\ud r\,\sigma(E).
  \]
  These polar rectangles form a generating family. Uniqueness of
  measures, first on finite annuli and then on their increasing union,
  identifies the polar pushforward with
  $r^{1-\alpha}\ud r\,\sigma(\ud\omega)$. Hence, for every
  nonnegative Borel function $f$, Tonelli gives
  \begin{equation}\label{E:planar-polar}
    \int f(x)\Gamma(\ud x)
    =\int_{\mathbb{S}^1}\int_0^\infty
    f(r\omega)r^{1-\alpha}\ud r\,\sigma(\ud\omega).
  \end{equation}
  Extended values are permitted. In particular, the formula applies
  to the nonnegative wave density used below.

  The endpoint polar measures are therefore
  $\ud r\,\sigma(\ud\omega)$ and
  $r^{1-\beta}\ud r\,\sigma(\ud\omega)$. Their radial integrals
  are finite near zero because $\beta<2$. Rapid decay of a Schwartz
  function controls their integrals at infinity, uniformly in direction.
  They are thus locally finite and tempered, and their stated scaling
  and evenness follow directly. If $\Gamma$ has a density, the two
  weighted measures also have locally integrable densities.

  To prove positive definiteness of the low endpoint, evenness gives
  the full-sphere line representation
  \begin{equation}\label{E:planar-line-mixture}
    \int f(x)\Gamma_{\rm low}(\ud x)
    =\frac12\int_{\mathbb{S}^1}\int_\R f(r\omega)\ud r
    \,\sigma(\ud\omega).
  \end{equation}
  It holds for nonnegative Borel $f$, and for Schwartz $f$ by absolute
  integrability. For a complex Schwartz function $\varphi$, put
  $\widetilde{\varphi}(x)=\overline{\varphi(-x)}$ and
  $\psi=\varphi*\widetilde{\varphi}$. Fubini in line and perpendicular
  coordinates gives
  \[
    \int_\R\psi(r\omega)\ud r
    =\int_{\omega^\perp}
    \left|\int_\R\varphi(s\omega+v)\ud s\right|^2\ud v\ge0.
  \]
  Schwartz decay justifies all these integrals, uniformly for unit
  $\omega$. Integrating against the finite measure $\sigma$
  in~\eqref{E:planar-line-mixture} proves
  $\Gamma_{\rm low}(\psi)\ge0$. The factor $1/2$ is essential because
  both orientations of each line occur on the full sphere.

  For the static endpoint put $\lambda=\beta-\alpha>0$. For $x\ne0$,
  \[
    |x|^{-\lambda}=\frac1{\Gamma(\lambda/2)}
    \int_0^\infty s^{\lambda/2-1}e^{-s|x|^2}\ud s.
  \]
  For each $s>0$, the measure $e^{-s|x|^2}\Gamma(\ud x)$ is
  positive definite. Indeed, the Fourier transform of
  $e^{-s|x|^2}\psi(x)$ is a positive Gaussian convolution of
  $|\widehat{\varphi}|^2$. Its pairing with the original positive
  spectral measure is nonnegative; Schwartz decay makes the pairing
  finite. Moreover,
  \[
    \int |\psi(x)||x|^{-\lambda}\Gamma(\ud x)
    =\int |\psi(x)|\Gamma_{\rm st}(\ud x)<\infty
  \]
  by the endpoint polar integrability just proved. Tonelli with
  $|\psi|$, followed by Fubini, justifies the Gaussian mixture against
  $\psi$. Thus, $\Gamma_{\rm st}(\psi)\ge0$. This uses smooth
  Gaussian multiplication of the original covariance and a convergent
  mixture, not a fractional power of that covariance.

  We next prove the convolution estimate directly. Use the planar wave
  density
  \[
    G_h(y)=\frac{\mathbf1_{\{|y|<h\}}}{2\pi\sqrt{h^2-|y|^2}},
  \]
  with value zero on the circle. Write
  $z=a\omega+b\omega^\perp$ for a unit perpendicular vector.
  If $|b|<h$, put $A=\sqrt{h^2-b^2}$. Translation along the line gives
  \[
    \int_\R G_h(z-r\omega)\ud r
    =\frac1{2\pi}\int_{-A}^A\frac{\ud s}{\sqrt{A^2-s^2}}
    =\frac12.
  \]
  If $|b|\ge h$, the open support is empty or tangent and the line
  integral is zero. The tangency point has zero line Lebesgue measure.
  Formula~\eqref{E:planar-line-mixture} and nonnegative Tonelli yield
  \[
    (G_h*\Gamma_{\rm low})(z)
    =\frac14\int_{\mathbb{S}^1}
    \mathbf1_{\{|z\cdot\omega^\perp|<h\}}\sigma(\ud\omega)
    \le\frac S4.
  \]
  Evenness, Gaussian convolution, and Tonelli now give
  \[
    \int G_{h,\eta}(\ud y)(p_\epsilon*\Gamma_{\rm low})(y+z)
    =\left[p_{\epsilon+\eta}*(G_h*\Gamma_{\rm low})\right](z)
    \le\frac S4.
  \]
  For $\eta=0$, use the unsmoothed Green measure. The Gaussian average
  has mass one, so the same constant works for all widths.

  Finally, the measure $p_\epsilon(z-x)\Gamma(\ud x)$ is finite.
  Gaussian decay and polynomial radial growth prove this directly;
  the endpoint convolutions are finite for the same reason. The identity
  \[
    \vartheta(\alpha-1)+(1-\vartheta)(\alpha-\beta)=0
  \]
  and absence of an origin atom permit H\"older's inequality on this
  unnormalized positive measure:
  \[
    \begin{split}
      (p_\epsilon*\Gamma)(z)
      &=\int |x|^{\vartheta(\alpha-1)}
      |x|^{(1-\vartheta)(\alpha-\beta)}
      p_\epsilon(z-x)\Gamma(\ud x)\\
      &\le(p_\epsilon*\Gamma_{\rm low})(z)^\vartheta
      (p_\epsilon*\Gamma_{\rm st})(z)^{1-\vartheta}.
    \end{split}
  \]
  This proves Lemma~\ref{L:planar-endpoints}.
\end{proof}

\begin{lemma}[General planar static coefficient input]
  \label{L:planar-static-input}
  Let $\Xi$ be a nonnegative, even, positive-definite tempered Radon
  measure on $\R^2$ with
  $\Xi(rA)=r^{2-\beta}\Xi(A)$, where $1<\beta<2$.
  For static Stratonovich noise with covariance $\Xi$, data $(1,0)$,
  and wave operator $\Delta$, let $U_n(t,x)$ be its ordered integrals,
  with $U_0=1$. There is a finite $C_{\rm st}=C_{\rm st}(\Xi)>0$ such
  that, for $k\ge1$ and $t>0$,
  \begin{equation}\label{E:planar-static-input}
    0\le B_k(t):=\sum_{m=0}^{2k}\E[U_m(t,x)U_{2k-m}(t,x)]
    \le\frac{C_{\rm st}^{\,k}t^{(4-\beta)k}}{(k!)^{3-\beta}}.
  \end{equation}
  Also $B_0=1$.
\end{lemma}

\begin{proof}
  If $\Xi=0$, then $U_n=0$ for $n\ge1$, and hence $B_k=0$ for
  $k\ge1$. The choices $C_{\rm st}=1$ and $B_0=1$ give the
  conclusion. We may therefore assume $\Xi\ne0$.

  Fourier uniqueness transfers covariance homogeneity to its positive
  spectral measure $\mu_\Xi$:
  $\mu_\Xi(rA)=r^\beta\mu_\Xi(A)$. Its unit-ball mass is finite
  because the spectral measure is Radon. On dyadic annuli this gives
  \[
    \int_{\R^2}(1+|\xi|^2)^{-1}\mu_\Xi(\ud\xi)
    \le\mu_\Xi(B_1)
    +2^\beta\mu_\Xi(B_1)\sum_{j\ge0}2^{j(\beta-2)}<\infty.
  \]
  This verifies the static Dalang condition without an angular density.
  From Chen--Hu~\cite{chen.hu:25:hyperbolic}*{Theorem~3.8(i), p.~2165},
  every fixed order exists in $L^2$. The product and Wick formulas
  (2.17) and (2.19), pp.~2151--2152, identify $B_k$ with all pairings
  and all splits of the total order, including self-pairs and empty
  chains. Thus, every coefficient is finite and nonnegative.

  We make explicit the finiteness input in the source's coefficient
  upper rate. The variational constant $\mathcal{M}$ in
  Chen--Hu~\cite{chen.hu:25:hyperbolic}*{p.~2147} is finite under the
  displayed Dalang condition. The source's invocation on p.~2175
  identifies $2^{\beta/(4-\beta)}\mathcal{M}$ with the square-root
  variational quantity covered by
  Chen~\cite{chen:23:exponential}*{Theorem~1.1, p.~2}.
  That theorem explicitly states finiteness under Dalang's condition
  and permits generalized nonnegative positive-definite covariances.
  We import that finiteness, not a sharp value or its proof.

  From Chen--Hu~\cite{chen.hu:25:hyperbolic}*{(5.1), p.~2173,
  and (5.15), p.~2180}, with $p=2$, it follows that
  \[
    B_k(t)=t^{(4-\beta)k}B_k(1),\qquad
    \limsup_{k\to\infty}\frac1k
    \log\left((k!)^{3-\beta}B_k(1)\right)<\infty.
  \]
  Choose a finite $R>0$ larger than this upper limit. For all sufficiently
  large $k$, $(k!)^{3-\beta}B_k(1)\le e^{Rk}$. Increase $e^R$ to a
  finite positive constant dominating the $k$th roots of the finitely
  many remaining quantities. The resulting $C_{\rm st}$
  proves~\eqref{E:planar-static-input} by homogeneity. These are the source's
  general homogeneous-covariance statements, not a Riesz specialization.
  This proves Lemma~\ref{L:planar-static-input}.
\end{proof}

\begin{proposition}[Planar homogeneous-covariance upper bound]
  \label{P:planar-upper}
  Let $d=2$, $0<\alpha_0<1$, and $1<\alpha<2-\alpha_0$.
  Suppose that the nonzero spatial covariance $\Gamma$ is nonnegative,
  even, positive definite and tempered,
  with~\eqref{E:planar-measure-scaling}. Choose and fix
  \[
    \alpha_0<\vartheta<2-\alpha,\qquad
    b_0=\alpha_0/\vartheta,\qquad
    \beta=(\alpha-\vartheta)/(1-\vartheta).
  \]
  Use the endpoint measures and physical angular mass $S$ of
  Lemma~\ref{L:planar-endpoints}. Put
  \[
    Z_{b_0}=\frac2{(1-b_0)(2-b_0)},\qquad
    C_{\rm low}=Z_{b_0}S,\qquad
    C=C_{\rm low}^{\,\vartheta}C_{\rm st}^{\,1-\vartheta},
    \qquad q=3-\alpha,
  \]
  where $C_{\rm st}$ is from Lemma~\ref{L:planar-static-input} for
  $\Xi=\Gamma_{\rm st}$. For the Stratonovich solution with data
  $(1,0)$, every $k\ge1$, $t>0$, and $x\in\R^2$ satisfy
  \begin{equation}\label{E:planar-upper}
    A_k(t)\le\frac{C^kt^{(4-\alpha-\alpha_0)k}}{(k!)^{3-\alpha}},
    \qquad
    \E u^2(t,x)\le\exp\left(qC^{1/q}t^\kappa\right).
  \end{equation}
  Constants may depend on the covariance, $\alpha_0$, and the fixed
  choice of $\vartheta$, but not on $k,t,x$.
\end{proposition}

\begin{proof}
  The parameter interval is nonempty. We have $0<\vartheta<1$,
  $0<b_0<1$, and
  \[
    \beta-\alpha=\frac{\vartheta(\alpha-1)}{1-\vartheta}>0,
    \qquad
    2-\beta=\frac{2-\alpha-\vartheta}{1-\vartheta}>0.
  \]
  Thus, Lemma~\ref{L:planar-endpoints} applies, with its interpolation
  weight equal to $\vartheta$. The original spectral measure has ball
  growth of exponent $\alpha$, by the same Fourier homogeneity argument
  used in Lemma~\ref{L:planar-static-input}. A dyadic sum gives the
  time-dependent spectral condition (1.5) of
  Chen~\cite{chen:author2026:time-dependency}*{p.~3}, since
  $\alpha<2-\alpha_0$. The low endpoint has spectral ball exponent
  one and $b_0<1$; the static endpoint has exponent $\beta<2$.
  Hence, the imported existence and fixed-order convergence statements
  are used within their stated parameter ranges.

  Fix $k\ge1$, $t>0$, and common noise widths $\epsilon,\delta>0$.
  The spatial covariance on each edge has width $2\epsilon$.
  No Green smoothing is needed in dimension two. On the disjoint union
  over $m=0,\ldots,2k$ and $D\in\Pi_k$, define the finite positive
  measure $\nu_{k,t}$ by
  \[
    \int\Phi\ud\nu_{k,t}
    :=\sum_{m=0}^{2k}\sum_{D\in\Pi_k}
    \int_{[0,t]^m_<\times[0,t]^{2k-m}_<}\ud s
    \int\Phi(m,D,s,y)\prod_{l=1}^{2k}G_{h_l}(\ud y_l).
  \]
  Here each chain has its own root, $h_l=s_l-s_{p(l)}$, the $y_l$
  are its increments, and $x_l$ is its cumulative position. A zero-length
  chain has empty products. At fixed times the spatial mass is
  $\prod_lh_l$; bounded time simplexes and finite discrete sums therefore
  make $\nu_{k,t}$ finite.

  For $J\in\{\Gamma,\Gamma_{\rm low},\Gamma_{\rm st}\}$ write
  \[
    H_J=\prod_{(i,j)\in D}(p_{2\epsilon}*J)(x_i-x_j),\qquad
    T_b=\prod_{(i,j)\in D}|s_i-s_j|^{-b}.
  \]
  The Gaussian covariance functions are bounded and continuous at each
  fixed width: their positive spectral measures are tempered, so Gaussian
  weighting makes them finite. Let $Q_\delta$ be the product of the
  target's truncated temporal covariances. The source regularization gives
  $0\le Q_\delta\le T_{\alpha_0}=T_{b_0}^{\,\vartheta}$.
  Lemma~\ref{L:planar-endpoints}, at the common edge width, gives
  $H_\Gamma\le H_{\Gamma_{\rm low}}^\vartheta
  H_{\Gamma_{\rm st}}^{1-\vartheta}$.

  The Gaussian moment formula and the two separate chain reversals in
  Proposition~\ref{P:mixed} identify the target regularized coefficient
  with $\int Q_\delta H_\Gamma\ud\nu_{k,t}$. Define
  \[
    L_k^\epsilon(t)=\int T_{b_0}H_{\Gamma_{\rm low}}\ud\nu_{k,t},
    \qquad
    B_k^\epsilon(t)=\int H_{\Gamma_{\rm st}}\ud\nu_{k,t}.
  \]
  The latter is the actual regularized static coefficient, including
  every split and pairing. The former is only a deterministic integral
  with raw temporal power. H\"older's inequality on the full common
  measure, including its discrete variables, yields
  \begin{equation}\label{E:planar-coefficient-holder}
    A_k^{\epsilon,\delta}(t)
    \le\left(L_k^\epsilon(t)\right)^\vartheta
    \left(B_k^\epsilon(t)\right)^{1-\vartheta}.
  \end{equation}
  No interpolation of different temporal cutoff kernels is asserted.

  We bound $L_k^\epsilon$ directly. Orient each pair by increasing block
  label. Descending integration of increments, as in
  Lemma~\ref{L:d1-spatial}, removes at a closing label its only remaining
  covariance factor. Lemma~\ref{L:planar-endpoints} bounds that integral
  by $S/4$. At an opening label its paired covariance has already been
  removed, so the Green integral is its mass $h_l$. The $k$ opening
  increments sum to at most $2t$. Thus, the common spatial bound is
  \[
    (S/4)^k\prod_{i\in O_D}h_i
    \le(S/2)^k\frac{t^k}{k^k}.
  \]
  The temporal pairing sum is invariant under every permutation of its
  $2k$ labels. Each pairing integrates on the full time cube to
  $(Z_{b_0}t^{2-b_0})^k$. Partitioning by the two within-block orders,
  and then summing all splits, gives
  \[
    \begin{split}
      \sum_{m=0}^{2k}\sum_{D\in\Pi_k}
      \int_{[0,t]^m_<\times[0,t]^{2k-m}_<}T_{b_0}\ud s
      &=\frac{(2k)!}{2^kk!}(Z_{b_0}t^{2-b_0})^k
      \sum_{m=0}^{2k}\frac1{m!(2k-m)!}\\
      &=\frac{2^k}{k!}(Z_{b_0}t^{2-b_0})^k.
    \end{split}
  \]
  Symmetry is used for the sum, not for an individual pairing. The
  integral is finite because $b_0<1$. Multiplication gives
  \begin{equation}\label{E:planar-low-coefficient}
    L_k^\epsilon(t)
    \le\frac{(Z_{b_0}S)^kt^{(3-b_0)k}}{k!\,k^k}
    \le\frac{C_{\rm low}^{\,k}t^{(3-b_0)k}}{(k!)^2}.
  \end{equation}
  At fixed $\epsilon$, bounded covariance and the finite common measure
  make $B_k^\epsilon$ finite as well. These observations also justify
  finiteness in the H\"older estimate.

  We next take fixed-order limits. The target orders converge in $L^2$
  by Chen~\cite{chen:author2026:time-dependency}*{Theorem~4.1}.
  For the static orders the corresponding input is
  Chen--Hu~\cite{chen.hu:25:hyperbolic}*{Theorem~3.8(i), p.~2165}.
  For products of two fixed orders, the inequality
  \[
    \|YZ-Y'Z'\|_1
    \le\|Y-Y'\|_2\|Z\|_2+\|Y'\|_2\|Z-Z'\|_2
  \]
  gives $L^1$ convergence. There are only $2k+1$ split products, and
  a zero-order factor stays equal to one. Hence,
  $A_k^{\epsilon,\delta}(t)\to A_k(t)$ and
  $B_k^\epsilon(t)\to B_k(t)$ as the noise widths vanish. The uniform
  deterministic estimate~\eqref{E:planar-low-coefficient} needs no
  stochastic limit interpretation. Passing to limits
  in~\eqref{E:planar-coefficient-holder} gives
  \[
    A_k(t)\le
    \left(\frac{C_{\rm low}^{\,k}t^{(3-b_0)k}}{(k!)^2}\right)^\vartheta
    B_k(t)^{1-\vartheta}.
  \]
  Only now apply Lemma~\ref{L:planar-static-input}. No monotonicity
  in the spatial width or uniform static smoothed asymptotic was used.

  The parameter identities
  $\alpha=\vartheta+(1-\vartheta)\beta$ and
  $\alpha_0=\vartheta b_0$ imply
  \[
    2\vartheta+(1-\vartheta)(3-\beta)=3-\alpha,\qquad
    \vartheta(3-b_0)+(1-\vartheta)(4-\beta)=4-\alpha-\alpha_0.
  \]
  This yields the coefficient bound in~\eqref{E:planar-upper} with
  $C=C_{\rm low}^\vartheta C_{\rm st}^{1-\vartheta}$.
  The constant coefficient is $A_0=1$. For $q=3-\alpha>1$ and $z\ge0$,
  \[
    \sum_{k\ge0}\frac{z^k}{(k!)^q}
    =\sum_{k\ge0}\left(\frac{z^{k/q}}{k!}\right)^q
    \le\left(\sum_{k\ge0}\frac{z^{k/q}}{k!}\right)^q
    =e^{qz^{1/q}}.
  \]
  Proposition~\ref{P:mixed} supplies the nonnegative full series for
  the target solution. Take $z=Ct^{4-\alpha-\alpha_0}$ and use
  $(4-\alpha-\alpha_0)/q=\kappa$. This proves
  Proposition~\ref{P:planar-upper}.
\end{proof}

\begin{corollary}[Full planar parameter range]\label{C:planar-range}
  In dimension two, under Hypothesis~\ref{H:homog} and the standing
  covariance positivity assumptions, the upper bound
  $\E u^2(t,x)\le\exp(C_2t^\kappa)$ holds throughout
  $0<\alpha<2-\alpha_0$, with a finite constant $C_2$ depending on
  $\alpha_0$ and the spatial covariance.
  A nonnegative generalized covariance is interpreted as the homogeneous
  measure and its Gaussian approximations described above.
\end{corollary}

\begin{proof}
  For $0<\alpha\le1$, apply Proposition~\ref{P:low-alpha-upper}.
  For $1<\alpha<2-\alpha_0$, apply Proposition~\ref{P:planar-upper}.
  If the covariance vanishes, $u=1$ and $C_2=0$ suffices. This proves
  Corollary~\ref{C:planar-range}.
\end{proof}

\begin{remark}[Planar calibrations and remaining goals]\label{R:planar-scope}
  The exponent tends to $(4-\alpha)/(3-\alpha)$ as
  $\alpha_0\downarrow0$, consistently with the static source, without
  an interchange of limits or uniformity claim for the constants.
  For the nonzero higher-exponent case, Cauchy--Schwarz gives
  $\E u(t,x)\le\exp(qC^{1/q}t^\kappa/2)$; the imported first-moment
  lower bound and Proposition~\ref{P:Jensen} imply
  $2L_*\le qC^{1/q}$. No sharp constant is identified. The existing
  comparisons with Skorohod exponents, in their stated Riesz subclasses,
  concern different solution concepts and are not inputs here. The planar
  endpoint construction
  is dimension-specific; no general higher-exponent result in dimension
  three follows from it. Thus, the full original homogeneous-covariance
  goal, the sharp constant, novelty and formal verification remain open.
\end{remark}

